\documentclass[10pt,a4paper]{amsart}
\usepackage[margin=19mm]{geometry}
\usepackage{amsmath,amssymb,mathtools}
\usepackage[scr=rsfs,cal=euler]{mathalfa}
\usepackage{xfrac}
\usepackage[unicode,hidelinks,bookmarks=false]{hyperref}
\newcommand{\ie}{i.e.\ }
\newcommand{\cf}{cf.\ }
\newcommand{\resp}{respectively\ }
\newcommand{\Subsection}{\subsection}
\providecommand{\nobreakdash}{\nobreak\hbox{-}}
\setlength{\emergencystretch}{2em}
\multlinegap0pt
\usepackage{tikz}
\usepackage{float}
\usepackage{amssymb}
\usepackage{enumerate}
\newtheorem{thm}{Theorem}
\newtheorem{lem}{Lemma}
\renewcommand{\O}{\mathcal{O}}
\DeclareMathOperator{\Spec}{Spec}
\newcommand{\cc}{\mathbb{C}}
\newcommand{\q}{\mathbb{Q}}
\title{Tautological Pushforwards of Hilbert Schemes of Points on Curves and Surfaces}
\author{Bochao Kong}
\date{Source: arXiv:2608.25447v1 (CC BY 4.0)}
\begin{document}
\maketitle

\noindent\emph{Source and licence.} Tautological Pushforwards of Hilbert Schemes of Points on Curves and Surfaces. Version arXiv:2608.25447v1 (CC BY 4.0), distributed by arXiv under the Creative Commons Attribution 4.0 International licence, http://creativecommons.org/licenses/by/4.0/, as recorded in the arXiv licence metadata for this exact version and shown on the abstract page https://arxiv.org/abs/2608.25447v1. Full licence notice: \texttt{licenses/arxiv-2608.25447-CC-BY-4.0.txt}.\par\smallskip

\noindent\textit{Editorial scope.} Counted unit: the article's thm thm:remaining-surface-codim-one-coefficients (source0.tex lines 468--485). The article's complete original proof of it is delivered with it (source0.tex lines 2903--3034, one \texttt{proof} environment of 132 lines, which the article places away from the statement). No mathematics is written, repaired or re-derived: the statement and the proof are the article's own lines, in the article's own notation, and no retained line is substituted or re-flowed.  Retained uncounted as the source-local context the proof needs: lines 408--413; lines 1175--1187; lines 2827--2838; lines 2848--2854; lines 2864--2877.  Nothing was omitted from any retained span, so the package carries no omission record.  No retained line contains a citation, so the wrapper carries no bibliography and no bibliography entry.  No retained line contains a figure environment, an image-inclusion command or drawing code, so no image is delivered and the package declares no asset dependency.  Editorial formatting only, in addition: an AMS \texttt{amsart} preamble that restates the article's own declarations and macros used by the retained lines, namely the environments \texttt{thm}, \texttt{lem} on separate counters, together with the standard packages those lines need.  The retained environments print in this document as Theorem 1, Lemma 1 in order of appearance, whereas the article prints the same items with the numbering of its own full document.  The complete native source of the article as distributed in the arXiv e-print for this exact version is bundled unchanged as \texttt{sources/208581/source0.tex}.
  \begin{equation}\label{surface-horizontal-recursion}
    A_{a,b,c}
    =
    -\frac1a
    \bigl(3\theta+a+b+2c-3\bigr)A_{a-1,b,c}.
  \end{equation}

\begin{lem}\label{pure-tangent-detector}
  Let $p_k(s,t)=s^k+t^k$. For every $k\ge2$, there is a smooth projective
  toric surface $S$, with isolated fixed points for its dense
  two-dimensional algebraic torus $T\cong(\mathbb G_m)^2$, such that
  \[
    (\pi_T)_{\star}p_k(T_{S_T/B_T})
    \in A^{k-2}(B_T)_\q=A_T^{k-2}(\Spec\cc)_\q
  \]
  is nonzero.
  Here $\pi_T:S_T\to B_T$ is the Borel mixed construction, and
  $p_k(T_{S_T/B_T})$ means evaluation of $p_k$ on the two Chern
  roots of the relative tangent bundle.
\end{lem}

\begin{align}
  [\zeta^e]\log Z^T_{S,L}(q,\zeta)
  &=
  \sum_{a+b+2c=e+2}
  A_{a,b,c}(q)(\pi_T)_{\star}(h^a t_1^b t_2^c) \notag\\
  &=
  \sum_{a+b+2c=e+2}
  A_{a,b,c}(q)
  \sum_{p\in S^T}
  \frac{h_p^a(x_p+y_p)^b(x_py_p)^c}{x_py_p},
  \label{compact-logarithmic-vertex-comparison}
\end{align}

\begin{equation}\label{local-edge-expansion}
  [\zeta]\log V(q,\zeta;0,x,y)
  =
  C_{011}^{\mathrm{loc}}(q)(x+y)
  +
  C_{030}^{\mathrm{loc}}(q)\frac{(x+y)^3}{xy}.
\end{equation}

  \begin{align}
    C_{030}^{\mathrm{loc}}(q)
    &=
    -\frac{v^2}{1-2v}
    \left(
      \frac12-12v+\frac{145}{2}v^2-147v^3+93v^4
    \right), \label{local-c030-closed-form}\\
    C_{011}^{\mathrm{loc}}(q)
    &=
    \frac{v^2}{1-2v}
    \left(
      1-16v+\frac{215}{4}v^2-\frac{129}{2}v^3+\frac{51}{2}v^4
    \right). \label{local-c011-closed-form}
  \end{align}

\begin{thm}\label{thm:remaining-surface-codim-one-coefficients}
  The remaining normalized codimension-one coefficients are
  \begin{align}
    A_{0,3,0}(v)
    &=
    -
    \frac{
      v^2\bigl(1-24v+145v^2-294v^3+186v^4\bigr)
    }
    {2(1-2v)}, \\
    A_{0,1,1}(v)
    &=
    \frac{
      v^2\bigl(102v^4-258v^3+215v^2-64v+4\bigr)
    }
    {4(1-2v)} .
  \end{align}
\end{thm}

\begin{proof}[Proof of Theorem
  \ref{thm:remaining-surface-codim-one-coefficients}]
  We begin with the tangent direction visible to compact toric localization.
  Take $S=\mathbb F_1$, with the tangent weights used in the proof of Lemma
  \ref{pure-tangent-detector}, and take $L=\O$ with the trivial
  linearization.  Then
  \[
    \sum_p(x_p+y_p)=0,\qquad
    \sum_p\frac{(x_p+y_p)^3}{x_py_p}=-2(X+2Y).
  \]
  Comparing \eqref{compact-logarithmic-vertex-comparison} with
  \eqref{local-edge-expansion} in degree $\zeta^1$ gives
  \[
    \bigl(C_{030}^{\mathrm{loc}}-A_{0,3,0}\bigr)(-2(X+2Y))=0.
  \]
  Hence
  \[
    C_{030}^{\mathrm{loc}}=A_{0,3,0}.
  \]

  To recover $A_{0,1,1}$, we must leave the slice $h=0$.  For a
  character $\eta$, the affine vertex satisfies
  \[
    V(q,\zeta;h+\eta,x,y)
    =
    V\left(\frac{q}{(1+\zeta\eta)^3},
    \frac{\zeta}{1+\zeta\eta};h,x,y\right).
  \]
  Indeed, for a rank $n$ bundle $E$,
  \[
    s_\zeta(E\otimes\eta)
    =(1+\zeta\eta)^{-n}s_{\zeta/(1+\zeta\eta)}(E),
  \]
  and the vertex has the normalizing factor $\zeta^{-2n}$. Differentiating the
  logarithm gives
  \[
    \partial_h\log V
    =-\zeta(3\theta+\zeta\partial_\zeta)\log V.
  \]
  Therefore \eqref{local-edge-expansion} implies
  \begin{equation}\label{local-horizontal-edge}
    [h\zeta^2]\log V(q,\zeta;h,x,y)
    =
    -(3\theta+1)
    \left(
      C_{011}^{\mathrm{loc}}(x+y)
      +
      C_{030}^{\mathrm{loc}}\frac{(x+y)^3}{xy}
    \right).
  \end{equation}

  Now let $S=\mathbb P^2$, with torus weights $0,X,Y$ on the coordinates,
  and let $L=\O_{\mathbb P^2}(-1)$ with its natural linearization. At the
  three fixed points,
  \[
    (h;x,y)=(0;X,Y),\quad (X;-X,Y-X),\quad (Y;-Y,X-Y).
  \]
  For $L^{\otimes m}$, extract the coefficient of $m\zeta^2$ in
  $\log Z^T_{S,L^{\otimes m}}$.  Coefficientwise in $q$ and $\zeta$,
  both sides of
  the compact comparison are polynomials in the integer $m$.  Since the
  identity holds for every $m\in\mathbb Z$, their coefficients in $m$ may
  be compared.  Put
  \[
    I_{011}=\sum_p h_p(x_p+y_p),\qquad
    I_{030}=\sum_p h_p\frac{(x_p+y_p)^3}{x_py_p}.
  \]
  Directly,
  \[
    I_{011}=-2(X^2-XY+Y^2),\qquad
    I_{030}=-9(X^2-XY+Y^2).
  \]
  The local expression from \eqref{local-horizontal-edge} is
  \[
    -(3\theta+1)
    \left(C_{011}^{\mathrm{loc}}I_{011}
    +C_{030}^{\mathrm{loc}}I_{030}\right).
  \]
  We next spell out the universal contribution.  At a fixed point $p$,
  replacing $L$ by $L^{\otimes m}$ replaces
  $h_p=c_1^T(L|_p)$ by $mh_p$.  Taking $e=2$ in
  \eqref{compact-logarithmic-vertex-comparison} therefore gives
  \[
    [\zeta^2]\log Z^T_{S,L^{\otimes m}}
    =
    \sum_{a+b+2c=4}m^aA_{a,b,c}
    \sum_p\frac{h_p^a(x_p+y_p)^b(x_py_p)^c}{x_py_p}.
  \]
  Hence the coefficient linear in $m$ comes exactly from $a=1$:
  \[
    [m\zeta^2]\log Z^T_{S,L^{\otimes m}}
    =
    \sum_{b+2c=3}A_{1,b,c}
    \sum_p\frac{h_p(x_p+y_p)^b(x_py_p)^c}{x_py_p}.
  \]
  The equation $b+2c=3$ has the two nonnegative solutions
  $(b,c)=(1,1)$ and $(3,0)$.  For these two solutions, the inner
  fixed-point sums are respectively $I_{011}$ and $I_{030}$.  Thus
  \[
    [m\zeta^2]\log Z^T_{S,L^{\otimes m}}
    =A_{1,1,1}I_{011}+A_{1,3,0}I_{030}.
  \]
  Since $a=1$ and $a+b+2c=4$, the horizontal recursion
  \eqref{surface-horizontal-recursion} specializes to
  \[
    A_{1,1,1}=-(3\theta+1)A_{0,1,1},\qquad
    A_{1,3,0}=-(3\theta+1)A_{0,3,0}.
  \]
  Thus the universal expression is
  \[
    -(3\theta+1)
    \left(A_{0,1,1}I_{011}+A_{0,3,0}I_{030}\right).
  \]
  Comparing the two expressions and using
  $C_{030}^{\mathrm{loc}}=A_{0,3,0}$, we obtain
  \[
    (3\theta+1)
    \bigl(C_{011}^{\mathrm{loc}}-A_{0,1,1}\bigr)I_{011}=0.
  \]
  Here $I_{011}=-2(X^2-XY+Y^2)\ne0$, and
  $A_T^*(\Spec\cc)_\q=\q[X,Y]$ is a domain.  Moreover,
  $3\theta+1$ is a $\q$-linear automorphism of $\q[[q]]$.
  It follows that
  \[
    C_{011}^{\mathrm{loc}}=A_{0,1,1}.
  \]

  Hence $A_{0,3,0}=C_{030}^{\mathrm{loc}}$ and
  $A_{0,1,1}=C_{011}^{\mathrm{loc}}$. The closed forms
  \eqref{local-c030-closed-form} and \eqref{local-c011-closed-form} are exactly
  the two formulas in the proposition.
\end{proof}

\end{document}
